\chapter{\label{chap:intro}Introdução}

A correção automática de afinação vocal tornou-se uma tecnologia amplamente utilizada na produção musical, tanto para ajustes sutis de interpretação quanto para efeitos estéticos característicos. Em diferentes contextos musicais, esse recurso assumiu papel relevante não apenas como ferramenta técnica de correção, mas também como elemento expressivo e criativo. Nesse cenário, o estudo de soluções de correção vocal mostra-se importante tanto sob a perspectiva prática quanto acadêmica, especialmente no que se refere aos fundamentos de áudio digital, à estimação da frequência fundamental e aos algoritmos responsáveis pela detecção de pitch \cite{antares_autotune_pro,antares_autotune_manual,decheveigne2002yin,mauch2014pyin}.

Nesse contexto, este trabalho apresenta o desenvolvimento de um protótipo experimental de correção automática de afinação vocal, construído com apoio de bibliotecas de software livre e com foco na investigação, implementação e comparação de algoritmos de detecção de pitch. Como etapa inicial, foi realizado um levantamento de requisitos com um usuário experiente no uso de ferramentas de correção vocal, a fim de identificar expectativas, dificuldades percebidas e funcionalidades desejadas em uma solução desse tipo. A partir disso, foram estudados os fundamentos teóricos envolvidos e comparados, em Python, os algoritmos de autocorrelação, YIN, pYIN e SWIPE', em termos de precisão, estabilidade, custo computacional, latência e adequação ao problema proposto \cite{decheveigne2002yin,mauch2014pyin,delacuadra2007pitch,camacho2008swipe}.

Com base nessa comparação, a abordagem de melhor equilíbrio entre os atributos analisados --- o pYIN --- foi utilizada como base para a evolução prática do protótipo em C++, com apoio do framework JUCE, empregando a técnica TD-PSOLA na etapa de correção. A escolha dessa linguagem e desse framework está relacionada à necessidade de maior controle sobre desempenho e aproximação com o desenvolvimento de aplicações e plug-ins de áudio em contexto real \cite{cppreference_cpp,juce_website,juce_github}. Dessa forma, o objetivo central do trabalho é compreender as diferenças entre os algoritmos estudados e, a partir dessa base técnica, desenvolver uma solução experimental de correção automática de afinação vocal.

Além desta introdução, o restante deste trabalho está organizado da seguinte forma: o Capítulo~\ref{chap:autotune_contexto} apresenta a correção automática de afinação vocal no contexto da produção musical, abordando suas aplicações, ferramentas comerciais, acessibilidade e integração com DAWs e plugins de áudio; o Capítulo~\ref{chap:problema} define o problema de pesquisa e explicita o desafio técnico de reconstruir um autotune; o Capítulo~\ref{chap:proposta} apresenta a proposta do trabalho, incluindo sua visão geral, o levantamento das necessidades do usuário, os requisitos, o escopo e as etapas de desenvolvimento do protótipo; o Capítulo~\ref{chap:conceitos_basicos} reúne os conceitos fundamentais necessários para a compreensão da proposta, como áudio digital, frequência fundamental, pitch, harmônicos e detecção de pitch; o Capítulo~\ref{chap:algoritmos_pitch} descreve os algoritmos de detecção de pitch considerados neste trabalho; o Capítulo~\ref{chap:estudo_comparativo} detalha o estudo comparativo desses algoritmos, com seus conjuntos de dados, métricas e protocolo experimental; o Capítulo~\ref{chap:resultados} apresenta e discute os resultados da comparação e a escolha do algoritmo principal; o Capítulo~\ref{chap:implementacao} descreve a implementação do protótipo em C++ com o framework JUCE; o Capítulo~\ref{chap:conclusao} reúne as conclusões e as contribuições do trabalho; e o Capítulo~\ref{chap:cronograma} apresenta o planejamento do desenvolvimento e a organização do trabalho em sprints.